% ============================================================================
% Doc. 265 -- Levels of Correspondence Between Physical Formalisms:
%             Blume-Maleev / FFGFT as a Worked Example
%
% Purpose: A worked example for the IPI structural comparison framework
%          (Stefaan Vossen, José Guevara RA-architecture, June 2026).
%          Shows three distinct levels of correspondence between two
%          independently derived formalisms, using the ferrotoroidic
%          memory paper (Qureshi et al. 2026) as a concrete test case.
%
% Status: Methodological document. No new FFGFT physics. The FFGFT
%         interpretations are from Doc. 264; this document adds the
%         meta-level classification of what type of correspondence each
%         item represents.
% ============================================================================

\section*{Purpose}

This document provides a worked example for the structural comparison of
physical formalisms, developed for the IPI group's model-ordering work.
The example uses two independently derived mathematical descriptions of the
same experimental system -- the ferrotoroidic quaternary memory crystal
LiNi$_{0.8}$Fe$_{0.2}$PO$_4$ (Qureshi et al.\ 2026) -- to illustrate that
``correspondence between formalisms'' is not a single category but a
stratified relation with at least three distinct levels.

The two descriptions are:
\begin{itemize}
  \item \textbf{Blume-Maleev polarization matrix formalism}: the standard
        condensed-matter description, complete and quantitatively accurate,
        developed entirely within established experimental physics.
  \item \textbf{FFGFT (T$^4$-torus geometry)}: the geometric field theory
        framework, applied to the same system from its own first principles,
        without fitting or borrowing from the condensed-matter description.
\end{itemize}

\noindent Neither framework was designed with the other in mind. The
correspondences that emerge are therefore structurally informative -- they
reveal what mathematical content is shared versus what is framework-specific.

% ============================================================================
\section{The Three Levels of Correspondence}
% ============================================================================

\begin{table}[htbp]
\centering
\caption{Classification of correspondence types between the
Blume-Maleev and FFGFT descriptions of the ferrotoroidic system.}
\label{tab:levels}
\begin{tabular}{p{2.5cm}p{5.5cm}p{5.5cm}}
\toprule
\textbf{Level} & \textbf{Definition} & \textbf{Diagnostic criterion} \\
\midrule
\textbf{1. Near-equivalence} &
  Both formalisms express the same mathematical object, derived
  independently from different starting points. &
  The abstract structure (group, operator, invariant) is identical;
  only the derivation path and physical interpretation differ. \\
\addlinespace
\textbf{2. Structural overlap} &
  The same qualitative structure appears in both, but the quantitative
  content is not shared. &
  One formalism can be mapped onto the other for the structural
  skeleton, but not for numerical predictions. \\
\addlinespace
\textbf{3. Complementary} &
  Each formalism explains something the other cannot. &
  Neither is a subset or extension of the other; the union covers
  more than either alone. \\
\bottomrule
\end{tabular}
\end{table}

% ============================================================================
\section{The comparison operator -- made explicit}
% ============================================================================

The three-level classification in Table~\ref{tab:ebenen} assigns each
correspondence to a level. Up to this point, however, it remained
\emph{implicit} \emph{by which criteria} this assignment is made -- which
operator decides whether a phenomenon falls on Level~1, 2 or~3. This
section exposes that operator, so that it becomes inspectable and
contestable rather than resting on intuition.

\subsection{The six comparison axes}

Two formalisms are compared along six independent axes. On each axis the
content is either \emph{shared} (both formalisms deliver the same) or
\emph{not shared}:

\begin{enumerate}
  \item \textbf{Shared invariants}: Do both produce the same conserved
        quantities / topological invariants?
        (e.g.\ the conserved winding / domain number)
  \item \textbf{Shared operators}: Does the same abstract operator
        structure appear? (e.g.\ the same group
        $\mathbb{Z}_2\times\mathbb{Z}_2$, the same transformation rule)
  \item \textbf{Shared mathematical structures}: Do both share the
        qualitative scaffold -- antisymmetric coupling, projection factor --
        even if the numerical values differ?
  \item \textbf{Shared predictions}: Do both deliver the same
        \emph{quantitative} value for the same experiment?
  \item \textbf{Shared physical mechanisms}: Do both attribute the
        phenomenon to the \emph{same} physical cause?
  \item \textbf{Shared hardware implications}: Do the same consequences
        for real devices / measurements follow from both?
\end{enumerate}

\subsection{How the axes define the levels}

The level assignment is not intuition but reads directly off the pattern
of shared axes:

\begin{table}[htbp]
\centering
\caption{The comparison operator: which axes are shared at which level.
$\checkmark$ = shared, $\times$ = not shared, $\sim$ = shared only
qualitatively.}
\label{tab:operator}
\begin{tabular}{p{5.5cm}ccc}
\toprule
\textbf{Comparison axis} & \textbf{Level 1} & \textbf{Level 2} & \textbf{Level 3} \\
 & Near-equiv. & Overlap & Complement. \\
\midrule
1. Invariants            & $\checkmark$ & $\sim$       & $\times$ \\
2. Operators             & $\checkmark$ & $\sim$       & $\times$ \\
3. Mathematical structure & $\checkmark$ & $\checkmark$ & $\times$ \\
4. Predictions (quantitative) & $\times$ & $\times$    & $\times$ \\
5. Physical mechanisms   & $\times$     & $\times$     & $\times$ \\
6. Hardware implications & $\times$     & $\times$     & $\times$ \\
\bottomrule
\end{tabular}
\end{table}

\noindent The assignment criterion is thus made explicit:
\begin{itemize}
  \item \textbf{Level 1 (near-equivalence):} Axes 1--3 fully shared --
        identical invariant, identical operator, identical structure --
        with different derivation paths. Axes 4--6 remain separate.
        (Examples: $\mathbb{Z}_2\times\mathbb{Z}_2$, $\mathcal{P}''=0$,
        topological charge.)
  \item \textbf{Level 2 (structural overlap):} Axis~3 shared
        (the same qualitative scaffold), axes 1--2 only partially, axes
        4--6 not. (Examples: antisymmetric exchange,
        decompactification projection.)
  \item \textbf{Level 3 (complementary):} None of axes 1--3 shared;
        each formalism covers a domain the other cannot reach.
        (Examples: geometric necessity of
        $\mathbb{Z}_2\times\mathbb{Z}_2$ in FFGFT; material-specific
        transition temperatures in Blume-Maleev.)
\end{itemize}

Notably, axes~4--6 (quantitative predictions, physical mechanisms,
hardware) are shared at \emph{no} level in this comparison. This is not
a defect of the classification but its main finding: Blume-Maleev and
FFGFT share mathematical \emph{structure} (axes 1--3) but not physical
\emph{substance} (axes 4--6). This is precisely what distinguishes this
correspondence from a genuine equivalence
(see the section on the Heisenberg/Schrödinger analogy).

\subsection{Residuals: limits of the operator}

The operator made explicit here is itself not fully mature. Three
residuals remain open and are declared here, not hidden:

\begin{itemize}
  \item \textbf{Axis completeness:} Are six axes sufficient, or are
        dimensions missing (e.g.\ shared boundary conditions, shared
        symmetry-breaking patterns)? The list is a working state, not a
        proven basis.
  \item \textbf{Axis weighting:} All six axes are weighted equally here.
        Whether a shared invariant (axis~1) should weigh more than a
        shared structure (axis~3) is undecided.
  \item \textbf{Axis independence:} The axes are treated as independent
        but may not be -- a shared operator structure (axis~2) often
        implies a shared invariant (axis~1). The correlation structure of
        the axes is not analysed.
\end{itemize}

These residuals do not change the level assignment in the present case
(the pattern in Table~\ref{tab:operator} is robust), but they mark where
the operator itself requires further examination. The value of making it
explicit lies precisely here: a previously implicit but load-bearing
criterion becomes an inspectable object -- inspectable rather than
persuasive.

% ============================================================================
\section{Level 1 -- Near-Equivalences}
% ============================================================================

\subsection{The $\mathbb{Z}_2\times\mathbb{Z}_2$ domain structure}

\textbf{In Blume-Maleev:} The four magnetic domains are related by two
symmetry operations of the crystal: time-inversion $1'$ (maps $a \leftrightarrow
b$, the ferrotoroidic flip) and the 2-fold screw axis $2_1\|y$ (maps
$'\leftrightarrow''$, the orientation flip). Together these generate a
$\mathbb{Z}_2\times\mathbb{Z}_2$ group structure over the four domains.
This is derived from the crystal symmetry group P$1121/a'$ and the
magnetic space group analysis.

\textbf{In FFGFT:} The T$^4$ torus has two independent binary degrees of
freedom for closed winding loops: the temporal winding direction ($\pm\tilde{T}$,
i.e.\ whether the loop winds clockwise or anti-clockwise in the compact time
direction) and the spatial winding parity ($\pm c$, the two possible
orientations of the winding plane in the observable 3D subspace). These
generate the same $\mathbb{Z}_2\times\mathbb{Z}_2$ group structure over the
four winding classes.

\textbf{Assessment:} The same abstract group $\mathbb{Z}_2\times\mathbb{Z}_2$
is derived from two completely different starting points -- crystal symmetry
on one side, T$^4$ compactification geometry on the other. This is a
near-equivalence: the mathematical object is the same, the derivation paths
are independent.

\textbf{What it does not mean:} This does not imply that the crystal lattice
\emph{is} a T$^4$ torus, or that the FFGFT parameters describe the specific
crystal. The shared group structure is a mathematical fact about how binary
symmetries combine -- it does not identify the underlying physical mechanism.

\subsection{The $\mathcal{P}'' = 0$ condition}

\textbf{In Blume-Maleev:} The general polarization transformation is
$\mathbf{P}_f = \mathcal{P}'\mathbf{P}_i + \mathcal{P}''$. For
LiNi$_{0.8}$Fe$_{0.2}$PO$_4$, $\mathcal{P}'' = 0$ because the magnetic
structure factors are purely imaginary: $\mathrm{Re}(NM_{\perp i}^*) = 0$
and $\mathrm{Im}(M_{\perp y}M_{\perp z}^*) = 0$. This follows from the
antiferromagnetic head-to-tail configuration (moments related by inversion
are antiparallel).

\textbf{In FFGFT:} The recursion operator $\mathcal{R}$ (Doc.~203) on the
T$^4$ mode space decomposes into a rotational part (coherent mode
transformation, winding number preserved) and a creation part (spontaneous
mode generation at T$^4$ boundary nodes, Doc.~204). In the static T$^4$
ground state -- no expansion, no horizon dynamics -- the creation part
vanishes: $\mathcal{P}'' = 0$.

\textbf{Assessment:} Both derive the same vanishing of the creation/annihilation
term from a ``no spontaneous generation'' condition -- in the crystal from the
antiferromagnetic symmetry, in FFGFT from the static geometry. The mathematical
structure is the same; the physical mechanism is different but yields the same
formal result.

\subsection{Topological charge conservation}

\textbf{In the paper:} The memory is non-volatile -- domain populations
remain stable at $E = H = 0$ after field-cooling. The paper identifies
this as a consequence of the energy barrier separating toroidal domains:
switching requires applying a conjugate field $E\times H$ above a threshold.

\textbf{In FFGFT:} A mode with non-zero winding number $n$ cannot decay to
$n = 0$ without crossing the minimal energy of the T$^4$ boundary condition
(Doc.~204, Fractal Boundary Condition). The winding number is a topological
invariant -- it is quantized and conserved unless a field large enough to
overcome the topological barrier is applied.

\textbf{Assessment:} Both describe the same mathematical phenomenon -- a
discrete topological invariant protected by an energy barrier -- though one
frames it in terms of magnetic anisotropy and the other in terms of T$^4$
compactification geometry. Near-equivalence at the level of the conservation
law; different mechanisms at the level of the physical origin.

% ============================================================================
\section{Level 2 -- Structural Overlap Without Equivalence}
% ============================================================================

\subsection{Antisymmetric exchange and the non-closure theorem}

\textbf{In Blume-Maleev:} The Dzyaloshinskii-Moriya (DM) interaction is
$\mathcal{H}^{\mathrm{DM}} = \mathbf{D}\cdot(\mathbf{S}_1\times\mathbf{S}_2)$,
odd under spin exchange (antisymmetric), with $\mathbf{D}_{34} = -\mathbf{D}_{12}$
from inversion symmetry.

\textbf{In FFGFT:} The Non-Closure Theorem (Doc.~193) states that cross-products
of mode labels from \emph{different} winding classes are topologically
inadmissible -- they would generate hybrid modes that violate the closed
winding structure of T$^4$. The antisymmetric cross-product $S_i\times S_j$
maps structurally onto this: it couples two spin components (two different
winding directions) in a way that is antisymmetric under interchange.

\textbf{Assessment:} The antisymmetric cross-product structure is shared.
But the DM vector components $D^b$ and $D^c$ are material-specific quantities
($|D^b| < |D^c|$, at $35°$ from the $c$-axis) that have no FFGFT counterpart.
FFGFT can say that the antisymmetric structure is topologically constrained;
it cannot predict the specific DM vector orientation or magnitude for this
crystal.

\subsection{Decompactification projection}

\textbf{In Blume-Maleev:} The factors $p_{fx}/p_{ix}$ in the general
polarization matrix (Eq.~9) measure the efficiency with which the measurement
apparatus resolves the polarization state -- how much of the abstract
mathematical object is made observable by the specific experimental setup.

\textbf{In FFGFT:} The decompactification projection (Layer~1 to Layer~2,
Doc.~001) measures the fraction of the unit-free geometric content that
becomes observable in the SI-projected description. The ratio
$T_{\mathrm{CMB}}/E_\xi = (16/9)\xi^2$ is a dimensionless Layer-1 prediction;
the absolute Kelvin value requires embedding in SI.

\textbf{Assessment:} The structure -- a projection factor that mediates
between the abstract mathematical content and the observable quantity -- is
the same in both. But the specific projection factors $p_{fi}$ are
apparatus-dependent, material-dependent numbers in the paper; in FFGFT the
projection is governed by $\xi = 4/30000$ and the SI embedding. No
quantitative mapping between the two is established.

% ============================================================================
\section{Level 3 -- Complementary (Non-Overlapping)}
% ============================================================================

\subsection{What FFGFT explains that Blume-Maleev does not}

The Blume-Maleev formalism derives the four-domain structure from the
observed crystal symmetry. It takes the symmetry group as given and
works out the consequences. It does not explain \emph{why} the relevant
symmetry group is $\mathbb{Z}_2\times\mathbb{Z}_2$ rather than some other
group -- that is determined by the specific material and is taken from
experiment.

FFGFT provides a geometric explanation: $\mathbb{Z}_2\times\mathbb{Z}_2$
is the \emph{necessary} structure for a compact T$^4$ torus with one
time-like and one space-like compact direction contributing to the observable
domain classification. The four-fold structure is not contingent on the
material; it is a consequence of the topology.

Similarly, the non-volatility of the memory is an empirical finding in
the paper. In FFGFT it is a structural necessity: topological charge
conservation is a consequence of the T$^4$ boundary condition, not
something that needs to be verified separately for each material.

\subsection{What Blume-Maleev provides that FFGFT does not}

The paper provides what FFGFT cannot: quantitative predictions for the
specific crystal. The transition temperatures $T_2 = 25$\,K and $T_4 = 21$\,K,
the domain population ratios (up to 90\% majority domain under optimal
field-cooling), the DM vector orientation ($35°$ from $c$-axis), the
polarization tensor elements $P_{xy}/P_{yx}$ and $P_{xz}/P_{zx}$ in
Table~1 -- all of these follow from the material parameters of
LiNi$_{0.8}$Fe$_{0.2}$PO$_4$ and are outside FFGFT's scope. FFGFT has
nothing to say about why the Ni/Fe ratio 0.8/0.2 produces four domains
rather than two, or why the transition temperature is 21\,K rather than
210\,K.

% ============================================================================
\section{The Analogy: Why Not Heisenberg vs.\ Schrödinger?}
% ============================================================================

A natural first analogy for two formalisms describing the same physics is
Heisenberg's matrix mechanics vs.\ Schrödinger's wave mechanics (1925--26),
which were subsequently proved to be fully equivalent (von Neumann 1932). But
this analogy is too strong for the present case.

Heisenberg and Schrödinger were describing the \emph{same physical objects}
(atoms, electrons) at the \emph{same scale} with the \emph{same predictive
scope}. Their equivalence meant that every prediction of one was also a
prediction of the other.

The Blume-Maleev / FFGFT correspondence is weaker: the two formalisms
share mathematical structure at Levels 1 and 2, but operate at different
scales (angstrom lattice vs.\ sub-Planck T$^4$), with different predictive
scope (material-specific quantitative vs.\ topological structural), and
different derivation bases (experimental symmetry analysis vs.\ geometric
first principles). There is no reason to expect full equivalence, and the
Level 3 analysis confirms it does not hold.

A more precise analogy is the relation between:
\begin{itemize}
  \item A symmetry-group classification (e.g.\ the Wigner classification of
        particles by irreducible representations of the Poincaré group): gives
        the topology, the quantum numbers, the conservation laws.
  \item A Lagrangian field theory for a specific particle (e.g.\ QED for the
        electron): gives all quantitative predictions for that specific particle.
\end{itemize}
The group classification is not a subset of QED, and QED is not a subset of
the Wigner classification. They are complementary: the group gives the
structural necessity (why the electron exists and has spin 1/2); the field
theory gives the dynamics (how it interacts). Neither alone is sufficient.

Blume-Maleev corresponds to the field theory (quantitative, material-specific,
dynamically complete). FFGFT corresponds to the group classification (structural
necessity, scale-independent, topologically explanatory). The shared content at
Levels 1 and 2 is the intersection -- the topological structure that both
descriptions must contain because it is a property of the physical system
independent of the formalism used to describe it.

% ============================================================================
\section{Implications for Model Comparison}
% ============================================================================

The three-level stratification has practical consequences for comparing
frameworks in the IPI group:

\begin{enumerate}
  \item \textbf{Level 1 correspondences are the most informative}: when two
        independently derived formalisms produce the same mathematical object
        (same group, same invariant, same conservation law), this is evidence
        that the object is a real structural feature of the physical system --
        not an artefact of the formalism. The $\mathbb{Z}_2\times\mathbb{Z}_2$
        structure appearing in both crystal symmetry analysis and T$^4$
        geometry suggests that this group is a genuine feature of toroidal
        topology, not just a convenient description.
  \item \textbf{Level 2 correspondences define the shared vocabulary}: the
        structural overlap tells us which concepts transfer between frameworks
        (antisymmetric coupling, projection factors) and which do not. This
        is useful for translation without overclaiming equivalence.
  \item \textbf{Level 3 defines the genuine boundary}: each framework's
        exclusive content marks where the other cannot substitute for it.
        Recognizing this boundary prevents both overclaiming (``my framework
        explains everything the other does'') and underclaiming (``the
        frameworks have nothing to do with each other'').
  \item \textbf{Full equivalence is the special case, not the default}: the
        Heisenberg/Schrödinger equivalence was a remarkable result, not a
        template. Most frameworks that share structural content are related by
        partial overlap or complementarity, not equivalence.
\end{enumerate}

\noindent\textbf{Reference.} The ferrotoroidic experiment used as the test
case: N.~Qureshi et al., \textit{Toroidicity as a route towards non-volatile
quaternary memory in antiferromagnets}, \textit{Nature Communications}
\textbf{17}, 4033 (2026). DOI:~10.1038/s41467-026-70767-8. The FFGFT
interpretation of all formulas in that paper is in Doc.~264 of this series.
